% GENERATED FILE. Edit canonical lessons, then run npm run generate:books.
% canonical-source-hash: fnv1a64:95eaca7b06e29a56
% canonical-lessons: FR-C303-tremblement-de-terre, FR-C303-avalanche, FR-C303-eboulement, FR-C303-effondrer, FR-C303-decombres

\chapter{Quand la terre tremble}
\label{ch:fr-quand-la-terre-tremble}
\hlchaptermodality{303}

\begin{chapteropening}
\textbf{By the end of this chapter:} \emph{I can talk about natural disasters: an earthquake, an avalanche, a rockfall, things caving in and the rubble.}

Complete the last lesson of chapter 303: I can talk about natural disasters: an earthquake, an avalanche, a rockfall, things caving in and the rubble.
\end{chapteropening}

\section[le tremblement de terre]{le tremblement de terre — the earthquake}

\label{lesson:FR-C303-tremblement-de-terre}



\begin{quote}
Before the new one: say the French for to explode, then the French for the spark.
\end{quote}

\subsection*{You'll want to know: le tremblement de terre}
\textbf{le tremblement de terre} — \textquotedblleft{}the earthquake\textquotedblright{}.

\begin{quote}
Le tremblement de terre a duré trente secondes.
\end{quote}

\begin{quote}
\emph{The earthquake lasted thirty seconds.}
\end{quote}

News reports also say \emph{le séisme}.

\begin{cousinweb}
\emph{tremblement} is built on \emph{trembler}, inherited from Vulgar Latin \emph{tremulāre}, built on Latin \emph{tremulus} \textquotedblleft{}trembling\textquotedblright{} (Lewis \& Short); \emph{terre} is inherited from Latin \emph{terra} \textquotedblleft{}earth\textquotedblright{} (EtymDB).
\end{cousinweb}

\begin{grammarlens}[title={What you've built}]
One more word for telling, explaining and describing at B1.
\end{grammarlens}

\subsection*{Your turn}
\begin{itemize}
\raggedright
  \item \emph{Say it:} \emph{le tremblement de terre}
  \item \emph{Say it:} \emph{le tremblement de terre}, once more
  \item \emph{Say it:} a sentence of your own with \emph{le tremblement de terre}
  \item \emph{Recall:} say the French for to explode, then the French for the spark, then say \emph{le tremblement de terre} again
\end{itemize}

\begin{tcolorbox}[breakable,colback=teal!4,colframe=teal!35!black,title={Before you move on}]
What does le tremblement de terre mean? (\textquotedblleft{}The earthquake\textquotedblright{}.) Say it once more.
\end{tcolorbox}
\section[l'avalanche]{l'avalanche — the avalanche}

\label{lesson:FR-C303-avalanche}



\begin{quote}
Before the new one: say the French for the spark, then the French for the earthquake.
\end{quote}

\subsection*{You'll want to know: l'avalanche}
\textbf{l'avalanche} — \textquotedblleft{}the avalanche\textquotedblright{}.

\begin{quote}
Une avalanche a fermé la route.
\end{quote}

\begin{quote}
\emph{An avalanche closed the road.}
\end{quote}

Figurative: \emph{une avalanche de lettres}.

\begin{cousinweb}
\emph{avalanche} was borrowed from Franco-Provençal \emph{avalançhe} (EtymDB), an Alpine word, shaped by \emph{avaler} \textquotedblleft{}to go down\textquotedblright{}.
\end{cousinweb}

\begin{grammarlens}[title={What you've built}]
One more word for telling, explaining and describing at B1.
\end{grammarlens}

\subsection*{Your turn}
\begin{itemize}
\raggedright
  \item \emph{Say it:} \emph{l'avalanche}
  \item \emph{Say it:} \emph{l'avalanche}, once more
  \item \emph{Say it:} a sentence of your own with \emph{l'avalanche}
  \item \emph{Recall:} say the French for the spark, then the French for the earthquake, then say \emph{l'avalanche} again
\end{itemize}

\begin{tcolorbox}[breakable,colback=teal!4,colframe=teal!35!black,title={Before you move on}]
What does l'avalanche mean? (\textquotedblleft{}The avalanche\textquotedblright{}.) Say it once more.
\end{tcolorbox}
\section[l'éboulement]{l'éboulement — the rockfall, the landslide}

\label{lesson:FR-C303-eboulement}



\begin{quote}
Before the new one: say the French for the earthquake, then the French for the avalanche.
\end{quote}

\subsection*{You'll want to know: l'éboulement}
\textbf{l'éboulement} — \textquotedblleft{}the rockfall, the landslide\textquotedblright{}.

\begin{quote}
Un éboulement a coupé la route de la montagne.
\end{quote}

\begin{quote}
\emph{A rockfall cut off the mountain road.}
\end{quote}

Signs in the mountains warn: \emph{chute de pierres}.

\begin{cousinweb}
\emph{éboulement} is built on \emph{ébouler} \textquotedblleft{}to crumble down\textquotedblright{}; the origin of that verb is not settled (EtymDB lists no parent).
\end{cousinweb}

\begin{grammarlens}[title={What you've built}]
One more word for telling, explaining and describing at B1.
\end{grammarlens}

\subsection*{Your turn}
\begin{itemize}
\raggedright
  \item \emph{Say it:} \emph{l'éboulement}
  \item \emph{Say it:} \emph{l'éboulement}, once more
  \item \emph{Say it:} a sentence of your own with \emph{l'éboulement}
  \item \emph{Recall:} say the French for the earthquake, then the French for the avalanche, then say \emph{l'éboulement} again
\end{itemize}

\begin{tcolorbox}[breakable,colback=teal!4,colframe=teal!35!black,title={Before you move on}]
What does l'éboulement mean? (\textquotedblleft{}The rockfall, the landslide\textquotedblright{}.) Say it once more.
\end{tcolorbox}
\section[effondrer]{effondrer — to collapse, to cave in (as s'effondrer)}

\label{lesson:FR-C303-effondrer}



\begin{quote}
Before the new one: say the French for the avalanche, then the French for the rockfall.
\end{quote}

\subsection*{You'll want to know: effondrer}
\textbf{effondrer} — \textquotedblleft{}to collapse, to cave in (as s'effondrer)\textquotedblright{}.

\begin{quote}
Le vieux pont s'est effondré.
\end{quote}

\begin{quote}
\emph{The old bridge collapsed.}
\end{quote}

Mostly reflexive. Also of a person: \emph{elle s'est effondrée en larmes} — she broke down in tears.

\begin{cousinweb}
\emph{effondrer} comes from Old French \emph{esfondrer} (EtymDB), built on \emph{fond} \textquotedblleft{}bottom\textquotedblright{}, from Latin \emph{fundus} \textquotedblleft{}bottom, ground\textquotedblright{} (Lewis \& Short).
\end{cousinweb}

\begin{grammarlens}[title={What you've built}]
One more word for telling, explaining and describing at B1.
\end{grammarlens}

\subsection*{Your turn}
\begin{itemize}
\raggedright
  \item \emph{Say it:} \emph{effondrer}
  \item \emph{Say it:} \emph{effondrer}, once more
  \item \emph{Say it:} a sentence of your own with \emph{effondrer}
  \item \emph{Recall:} say the French for the avalanche, then the French for the rockfall, then say \emph{effondrer} again
\end{itemize}

\begin{tcolorbox}[breakable,colback=teal!4,colframe=teal!35!black,title={Before you move on}]
What does effondrer mean? (\textquotedblleft{}To collapse, to cave in (as s'effondrer)\textquotedblright{}.) Say it once more.
\end{tcolorbox}
\section[les décombres]{les décombres — the rubble, the debris}

\label{lesson:FR-C303-decombres}



\begin{quote}
Before the new one: say the French for the rockfall, then the French for to collapse.
\end{quote}

\subsection*{You'll want to know: les décombres}
\textbf{les décombres} — \textquotedblleft{}the rubble, the debris\textquotedblright{}.

\begin{quote}
Les secours cherchent dans les décombres.
\end{quote}

\begin{quote}
\emph{The rescue teams are searching in the rubble.}
\end{quote}

Plural only.

\begin{cousinweb}
\emph{décombres} is built on Old French \emph{descombrer} \textquotedblleft{}to clear away\textquotedblright{}, the opposite of \emph{encombrer} \textquotedblleft{}to block\textquotedblright{}; the root is usually traced to a Gaulish word for a blockage, and the origin is not settled (EtymDB lists no parent).
\end{cousinweb}

\begin{grammarlens}[title={What you've built}]
One more word for telling, explaining and describing at B1.
\end{grammarlens}

\subsection*{Your turn}
\begin{itemize}
\raggedright
  \item \emph{Say it:} \emph{les décombres}
  \item \emph{Say it:} \emph{les décombres}, once more
  \item \emph{Say it:} a sentence of your own with \emph{les décombres}
  \item \emph{Recall:} say the French for the rockfall, then the French for to collapse, then say \emph{les décombres} again
\end{itemize}

\begin{tcolorbox}[breakable,colback=teal!4,colframe=teal!35!black,title={Before you move on}]
What does les décombres mean? (\textquotedblleft{}The rubble, the debris\textquotedblright{}.) Say it once more.
\end{tcolorbox}
